\documentclass[12pt]{article}

\usepackage[portuguese]{babel}
\usepackage[dvipsnames]{xcolor}
\usepackage{graphicx}
\usepackage{float}
\usepackage{parskip}
\usepackage{url}
\usepackage{geometry}
\usepackage{setspace}
\usepackage{classicthesis}

\geometry{
    left=112pt,
    right=112pt,
    bottom=72pt,
    top=72pt
}

\graphicspath{{./gfx/}}

\pagestyle{empty}

\title{
    \vspace*{-48pt}
    \spacedlowsmallcaps{\normalsize Resumo de Atividades} \\[4pt]
    \LARGE Mapeamento da Cobertura Hospitalar no Brasil para Doenças Tempo-Sensíveis
}
\author{
    \textit{Bolsista:} Bruno Samuel Ardenghi Gonçalves \\
    \textit{Orientador:} Prof. Dr. Márcio Dorn
}
\date{\normalsize Junho de 2025}

\begin{document}

\begin{figure}
    \vspace*{-36pt}
    \centering
    
\includegraphics[width=48pt]{ufrgs-sedetec.pdf}
\end{figure}

\maketitle
\thispagestyle{empty}

\setstretch{1.05}

\begin{figure}[H]
    \centering
    \includegraphics[width=\linewidth]{02.png}
\end{figure}

\section*{Introdução}

O \textit{MapStroke} (\href{https://sbcb.inf.ufrgs.br/mapstroke}{sbcb.inf.ufrgs.br/mapstroke}) é uma plataforma web criada para combater a desigualdade no acesso ao tratamento de Acidente Vascular Cerebral (AVC). A ferramenta auxilia gestores de saúde a identificar áreas carentes e alocar recursos de forma otimizada na instalação de novos centros de atendimento para doenças tempo-sensíveis.

A aplicação utiliza uma interface moderna e fluida, desenvolvida com \textit{SvelteKit}, \textit{TypeScript} e \textit{Tailwind CSS}, e emprega a biblioteca \textit{Leaflet} para a renderização de mapas interativos.

\section*{Funcionalidades Principais}

A funcionalidade central da plataforma é a ferramenta de visualização de cobertura por isócronas, que permite analisar a assistência hospitalar de uma região. Nela, o usuário define uma localização (país, estado, cidade ou hospital) e um tempo de deslocamento para que o sistema exiba no mapa os hospitais e a área total alcançável, juntamente com estatísticas de população e área coberta.

O acesso além dessa página é controlado por um sistema de solicitações. Interessados submetem um pedido de registro com dados profissionais que é avaliado por um administrador. Uma vez aprovado, o usuário obtém acesso às funcionalidades do papel a ele atribuído.

\section*{Painel Administrativo}

O painel administrativo centraliza um conjunto de ferramentas para a gestão da plataforma. A principal delas é o gerenciamento completo da base de hospitais. Ao criar um hospital ou modificar um existente (assistido de um mapa interativo), a isócrona correspondente à localização escolhida é imediatamente calculada e integrada à base de dados.

O painel também permite a aprovação ou rejeição de pedidos de acesso de usuários. Complementarmente, foi implementada a estrutura de um sistema de permissões baseado em papéis, que permitirá a configuração granular do que cada tipo de usuário pode fazer.

Por fim, a área administrativa permite a edição do conteúdo da página "Sobre" através de um editor Markdown com pré-visualização em tempo real.

\section*{Conclusão e Próximos Passos}

A plataforma MapStroke já possui suas funcionalidades centrais implementadas, incluindo a visualização de isócronas, um sistema de autenticação e um painel de administração robusto, estabelecendo uma base sólida para o projeto.

Os próximos passos do desenvolvimento se concentrarão em três frentes principais: o desenvolvimento de ferramentas analíticas para identificar automaticamente áreas críticas que necessitam de intervenção; a flexibilização da plataforma para permitir a análise de outras doenças e a incorporação de dados socioeconômicos; e a finalização do sistema de permissões, incluindo escopo de acesso por localização para adequar a plataforma a diferentes cenários de uso.

\end{document}
